The story begins in ancient Greece,
where discrete approximations were employed to obtain continuous areas through the method of exhaustion.
In particular, Archimedes (287--212 BC) evaluated sums of squares while determining
the area enclosed by his spiral~\cite{pengelley2002bridge},
leading to calculations analogous to the modern identity
\begin{align*}
  1^2 + 2^2 + 3^2 + \cdots + n^2 = \frac{n(n+1)(2n+1)}{6}.
\end{align*}

Moving to early modern Europe, Pierre de Fermat (1601--1665)
described the problem of finding formulas for sums of powers as
\textit{``perhaps the most beautiful problem of all arithmetic''}
in his correspondence with Blaise Pascal (1623--1662).
Fermat's contribution to the subject was substantial:
he revived the study of figurate numbers and investigated
their arithmetic properties,
including what is now known as Fermat's polygonal number theorem.
Applications of polygonal numbers to sums of powers
are discussed extensively in~\cite[(2020)]{Marko2020}.

Meanwhile, Johann Faulhaber (1580--1635) investigated sums of powers
in his celebrated work \textit{Academia Algebrae}~\cite[(1631)]{faulhaberacademia}.
His achievements were extraordinary,
as he obtained formulas up to
$\KnuthRFoldSum{1}{n}{25}=1^{25}+2^{25}+\cdots+n^{25}$.
Faulhaber explicitly tabulated formulas through the seventeenth power
and provided sufficient information for determining higher cases,
including the twenty-fifth power; the latter reconstruction was later carried out by Knuth.
More importantly, Faulhaber discovered deep structural patterns,
particularly the representation of odd-power sums as polynomials
in triangular numbers.
What remained missing was a completely general symbolic formula.

A major breakthrough was achieved by Jacob Bernoulli (1655--1705)
in his posthumously published work
\textit{Ars Conjectandi}~\cite[(1713)]{bernoulli1713ars}.
There, Bernoulli introduced the numbers that now bear his name,
obtaining a general formula for arbitrary power sums
\begin{align*}
  \sum_{k=1}^{n} k^p
  =
  \frac{1}{p+1}
  \sum_{r=0}^{p}
  \binom{p+1}{r}
  B_r n^{p+1-r}.
\end{align*}
Here it is assumed that $B_1=+\frac{1}{2}$.
The choice of convention for $B_1$ has a long historical background;
see~\cite{luschny2021manifesto,luschny2020introduction}
for a detailed discussion.
Bernoulli's formula supplied the missing symbolical framework
for Faulhaber's observations, thus the theory were complete.
Although the expression above is commonly known as Faulhaber's formula,
its modern Bernoulli-number formulation originates from Bernoulli's work.

During the eighteenth and nineteenth centuries,
the theory was further developed by Euler, Stirling, Jacobi,
and many others through connections with finite differences,
interpolation, generating functions, and special number sequences.

In 1993, Donald Knuth revisited the theory of sums of powers
and opened a new direction of development.
Building upon ideas of Carl Gustav Jacobi (1804--1851)
and the work of Gessel and Viennot~\cite[(1989)]{gessel1989determinants},
Knuth reexamined Faulhaber's results through the lens of modern
enumerative combinatorics.
Remarkably, Bernoulli numbers are not central to Knuth's approach.
Instead, he derived formulas for power sums using
special number sequences, including Stirling and central factorial numbers,
together with binomial coefficients.
Knuth also introduced the convenient notation
$\KnuthRFoldSum{r}{n}{m}$ for $r$-fold sums of powers,
which will be used throughout this paper.

In this manuscript, we show the role of Newton's interpolation formula
in the derivation of identities for sums of powers.
Our main result is an algorithm for obtaining closed-form expressions
for power sums by combining variants of Newton's interpolation formula
with binomial identities from the hockey-stick family.
The results demonstrate that infinitely many distinct closed-form representations
of sums of powers can be constructed in this manner.
